\section{Background and Model}
\label{sec:background}

This section reviews level-synchronous and direction-optimizing BFS, defines the progress guarantees we use, relates a wait-free search to the Do-All problems of asynchronous parallel computation, and states the model under which our method is wait-free. Throughout the paper, we say \emph{vertex} for an element of a graph and \emph{node} for an element of a linked list.

\subsection{Breadth-First Search}
\label{sec:bg-bfs}

We treat every graph as undirected. A graph $G = (V, E)$ has $n$ vertices, numbered $0$ to $n-1$, and we store each of its edges $\{u, v\}$ as the two directed edges $(u, v)$ and $(v, u)$; $m$ counts these directed edges. The graph is kept in compressed sparse row (CSR) form: the neighbors of a vertex $u$ are $\mathit{adj}[\mathit{offsets}[u]]$ up to, not including, $\mathit{adj}[\mathit{offsets}[u+1]]$, and the \emph{degree} of $u$ is $\mathit{offsets}[u+1] - \mathit{offsets}[u]$.

A BFS from a source vertex $s$ computes the distance $\delta(v)$ of every vertex $v$: the number of edges on a shortest path from $s$ to $v$, or $-1$ if there is no such path. Level $d$ is the set $L_d = \{v : \delta(v) = d\}$, so $L_0 = \{s\}$, and $D$ is the largest distance. The search keeps the distances it has found in an array $\mathit{dist}$, which holds $-1$ for every undiscovered vertex. A vertex is \emph{discovered} when its distance is set, and \emph{expanded} when its neighbors are examined; the level being expanded is the \emph{frontier}.

A top-down step expands the frontier $L_d$: it examines every neighbor $v$ of every vertex in $L_d$, and discovers each $v$ that is still undiscovered, setting $\mathit{dist}[v] = d+1$ and adding $v$ to $L_{d+1}$. A level-synchronous parallel BFS~\cite{bader_bfs, charles_bfs} divides each frontier among the threads; when several threads find the same undiscovered vertex, a compare-and-swap (CAS) on its distance lets only one of them add it to the next level. The step ends at a barrier, since no thread may expand $L_{d+1}$ before all of it has been discovered: otherwise a thread could reach a vertex of $L_{d+1}$ that is still undiscovered from another vertex of $L_{d+1}$, and give it the distance $d+2$.

A bottom-up step~\cite{beamer_dobfs} works the other way around: every undiscovered vertex $v$ scans its own neighbors for one in the frontier, and if it finds one, it sets $\mathit{dist}[v] = d+1$, adds $v$ to $L_{d+1}$, and stops scanning. Only the thread that scans $v$ writes its distance, so the step needs no CAS, but it needs a fast test of whether a vertex is in the frontier, such as the bitmap of GAPBS, and it too ends at a barrier. When the frontier is large, a bottom-up step examines far fewer edges than a top-down step: most undiscovered vertices find a frontier neighbor among their first few neighbors, while a top-down step examines every edge of the frontier, most of which lead to vertices that are already discovered.

A direction-optimizing BFS chooses the direction of every step. The rule of Beamer et al., as GAPBS implements it, keeps two counts: $m_f$, the number of edges of the frontier, which is the sum of the degrees of its vertices, and $m_u$, an estimate of the edges not yet examined, which starts at $m$ and drops by $m_f$ with every top-down step. A search going top-down switches to bottom-up once $m_f > m_u / \alpha$. A search going bottom-up continues while the frontier does not shrink, or holds more than $n / \beta$ vertices, and otherwise tests $m_f > m_u / \alpha$ with $m_f = 1$, which switches it back to top-down unless $m_u < \alpha$. GAPBS sets $\alpha = 15$ and $\beta = 18$; our method uses the same rule, which Algorithm~\ref{alg:wf-do} states in full (\textsc{Rule}), and our implementation the same values.

\subsection{Progress Guarantees}
\label{sec:bg-progress}

Progress guarantees are properties of methods~\cite{aomp}. A method is \emph{blocking} if a delay of one thread can keep the calls of other threads from finishing, as when a call waits for a lock that a delayed thread holds, or at a barrier that a delayed thread has not reached; a level-synchronous BFS is blocking, since a thread delayed before the barrier of a step delays every other thread there, and one that stops for good keeps them waiting forever. A method is \emph{lock-free} if it guarantees that infinitely often some call finishes in a finite number of steps: the system as a whole makes progress, although a particular call may never finish. A method is \emph{wait-free} if every call finishes in a finite number of the calling thread's own steps, whatever the other threads do~\cite{herlihy_waitfree}, and \emph{bounded wait-free} if one bound on that number holds in every execution~\cite{herlihy_waitfree, aomp}. The bound may depend on fixed parameters, such as the number of threads, but not on the schedule; in a method that is wait-free but not bounded, a call can take more steps the more other threads have done, as long as it finishes. A call of a wait-free method therefore never waits for another thread, and even a thread that stops for good delays no other. The usual way to make a method wait-free is \emph{helping}: instead of waiting for a delayed thread to finish its part of the work, a thread finishes that part itself~\cite{aomp, kogan_fastwf}.

\subsection{Wait-Free Algorithms and Do-All}
\label{sec:bg-doall}

Wait-freedom is hard to achieve when the threads hold inputs that the others do not know: the operations of a concurrent object, or the proposals of a distributed task such as consensus, which no wait-free protocol can solve with reads and writes~\cite{flp, herlihy_waitfree, herlihy_shavit_topology}. A search is different. Every thread knows the whole input, the graph and the source, and the output is a function of it. Wait-freedom alone is then easy to have: each thread runs a sequential BFS by itself and writes the distances, all threads writing the same values, so that every call returns after $O(n+m)$ of its own steps, whatever the other threads do. This trivial method does $T$ times the work of a sequential search, and is no faster than one. The question is how much less work, and how much more speed, wait-freedom allows.

That question has been studied for a basic problem of the same kind. In \emph{Write-All}, $P$ asynchronous or crash-prone processors must write 1 into every cell of an array of $N$ cells~\cite{kanellakis_robust}; \emph{Do-All} generalizes the cells to $N$ tasks~\cite{georgiou_doall}, and \emph{Certified Write-All} (CWA) adds a flag that may be set only once every cell is written, which abstracts telling that a set of tasks is complete~\cite{martel_cwa}. The measure is \emph{work}, the total number of steps of all processors~\cite{martel_apram}, by which Anderson and Woll also measured wait-free union-find~\cite{anderson_woll_uf}. Letting every processor do every task takes work $\Theta(PN)$; deterministic and randomized algorithms do far better, and lower bounds are known for several models~\cite{kanellakis_restartable, martel_cwa, anderson_woll_cwa, chlebus_writeall, kowalski_writeall}. Anderson and Woll let each processor visit the cells in a permutation of its own, and bound the redundant work by a property of the permutations that they call their \emph{contention}~\cite{anderson_woll_cwa}; Section~\ref{sec:pf-work} uses their argument. A CWA per step runs a synchronous PRAM program on asynchronous processors~\cite{martel_apram, kanellakis_book}, as asynchronous PRAM models, which make the speeds of processors part of the model, call for~\cite{cole_apram, gibbons_apram, nishimura_async}, and iterating Do-All in this way has been analyzed as a problem of its own~\cite{georgiou_iterative}.

A level of a BFS is an instance of Do-All: its tasks are the vertices of the level, each weighted by its degree, and the next level may start only once all of them are done. Run with one CWA per level over the $n$ vertices, a search does work at least $n$ in every level, $\Omega(Dn)$ in all, however few vertices the levels hold; on the road network of our experiments, with 6,262 levels and 24 million vertices, that is $1.5 \times 10^{11}$, against 58 million directed edges. This cost belongs to that mapping, not to every construction built on CWA. Three things set a level apart from the instances that this literature studies. Its tasks are discovered during the level before, by the same threads, so there is no array of tasks when the level begins. Its tasks have the degrees of the graph as their weights, as skewed as the graph's degrees are. And with direction optimization, the threads must also agree on how to expand it. Section~\ref{sec:wf-topdown} shows how our method handles the first, Section~\ref{sec:pf-work} what the second costs, and Section~\ref{sec:dir-opt} how it handles the third.

\subsection{System Model}
\label{sec:bg-model}

There are $T \geq 1$ threads, numbered $0$ to $T-1$, which communicate through shared memory. The number of threads is fixed before a search, and every thread knows it and its own number. Threads are asynchronous: any thread can be delayed for any length of time between any two of its steps, or stop for good, and we assume nothing about their relative speeds. The graph has $n \geq 1$ vertices, the source of a search is one of them, and the graph does not change during a search.

The shared memory is made of words of up to eight bytes, each aligned to its size, and every value that the method keeps in a word, such as a vertex, a distance, a position in a bucket or a count of edges, fits in it. Every shared variable of the method, such as a distance, a flag, the size of a bucket or a pointer, is a word, and a plan takes a few words. A thread accesses a word with a load, a store or a CAS, each of which is atomic, and in our algorithms every read of a shared variable is one load, every write one store, and every call of \textsc{CAS} one CAS. Memory follows total store order (TSO), the memory model of x86 processors~\cite{x86tso, intelmanual, amdmanual}. The stores of a thread become visible to all other threads at the same time, and in the order in which the thread made them, and the loads of a thread take effect in program order. The one reordering that TSO allows is that a load can take effect before an earlier store of the same thread to another location becomes visible, while the store still waits in the thread's store buffer. A CAS waits until every earlier store of its thread is visible, and then reads and writes its word in one atomic step; a \emph{fence} only waits. Both therefore order the accesses of their thread around them. Our method relies on the program order that TSO keeps between two accesses of a thread, a store and a later store, a load and a later load, or a load and a later store, and uses a CAS, or in the method of Section~\ref{sec:wf-topdown} one fence per level, where it needs a store to be visible before a later load. Each load, store, CAS or fence is one step of its thread. The method and its proofs use these operations only through this model, and Section~\ref{sec:impl} describes how our implementation realizes them.

Each thread allocates the memory it needs during a search from an arena of its own, and we assume that it can allocate memory in a bounded number of its own steps, and draw a random number from a generator of its own in a constant number. The bounds of Section~\ref{sec:proof} count an allocation and a random draw as a constant number of steps each; an allocator whose cost grows with $n$ or $T$ would leave the method wait-free, but would add its cost to these bounds. As wait-freedom usually is, our guarantee is stated for this model of primitive operations, not for the allocator, the operating system or the hardware below them. A method that may need unbounded memory loses wait-freedom when memory runs out~\cite{nikolaev_wcq}; a search, by contrast, needs a bounded amount of memory for a fixed graph and number of threads (Section~\ref{sec:mem-memory}), so an implementation can meet the assumption by reserving that memory before the search starts (Section~\ref{sec:impl}).

\subsection{The BFS Method}
\label{sec:bg-method}

Every thread calls the method \textsc{BFS}$(s)$ once per search, with the same source $s$. Before the first call, $\mathit{dist}[s] = 0$ and $\mathit{dist}[v] = -1$ for every other vertex $v$, the method's other arrays hold their initial values, and the first level holds the source; Sections~\ref{sec:td-structure} and~\ref{sec:dir-opt} give this initial state in full, and Section~\ref{sec:mem-reuse} that of the variant of Section~\ref{sec:variants}. Setting up this state, and resetting it between searches, is not part of the method.

When any call returns, $\mathit{dist}[v] = \delta(v)$ for every vertex $v$, and no step that any thread takes afterwards changes $\mathit{dist}$. A search is therefore complete when its first call returns, and that is when we stop timing it (Section~\ref{sec:exp-setup}). From then on, its result can be used while delayed calls may still be running, since they can only write the distances that are already there, and a thread whose call has returned is free to do other work, or to join another search that uses separate state. What waits for the last call is only the reuse or reclamation of the state of the search, which a delayed call may still access: the method guarantees the progress of its calls, not the reclamation of the state that they use, as nonblocking methods generally do, and wait-free reclamation is a subject of its own~\cite{nikolaev_crystalline}. A pool of threads can keep several copies of the state, so that a thread that stops for good withholds only one. Section~\ref{sec:proof} proves these properties, and that every call returns after a bounded number of the calling thread's own steps.
